% !TeX root = ../../Book1.tex
% 纲要标题：### 6.3 群扩张
% 纲要位置：计划纲要.md，第 286 行。

\section{群扩张}
\label{b1:sec:0603}

% 本节正文按下列原纲要要点撰写。

% 纲要内容（仅作写作计划，尚非教材正文）：
%
% * 短正合群列与分裂扩张
% * 截面是集合映射还是群同态
% * 扩张数据与半直积的关系
%

\subsection{短正合群列与分裂扩张}
\label{b1:subsec:0603:01}

商群把正规子群压缩为单位；群扩张提出相反的问题：已知正规子群和商群，能恢复哪些群？直积和半直积给出两种答案，但并非所有扩张都属于这两类。

\begin{definition}\label{b1:def:group-extension}
一列群同态
\[
1\longrightarrow N\xrightarrow{\iota}G
\xrightarrow{\pi}H\longrightarrow1
\]
称为\term[duan-zhenghe-qunlie]{短正合群列}[short exact sequence of groups]，如果 \(\iota\) 单射、\(\pi\) 满射，且 \(\operatorname{im}\iota=\ker\pi\)。它也称为 \(H\) 由 \(N\) 的一个\term[qun-kuozhang]{群扩张}[group extension]。若存在群同态 \(s:H\rightarrow G\) 满足 \(\pi s=\operatorname{id}_H\)，称扩张\term[fenlie-kuozhang]{分裂}[split extension]。
\end{definition}
这里的 \(1\) 表示平凡群。“正合”在每个位置要求前一个映射的像等于后一个映射的核，因此两端的条件正好是单射与满射。通过 \(\iota\)，可将 \(N\) 识别为 \(G\) 的正规子群；此时 \(G/N\cong H\)。但扩张还记住了指定映射 \(\iota,\pi\)，比只写一个抽象同构 \(G/N\cong H\) 包含更多信息。

\begin{theorem}\label{b1:thm:split-extension-semidirect}
短正合群列分裂，当且仅当 \(G\) 含有子群 \(L\)，使 \(\pi|_L:L\rightarrow H\) 为同构。存在这样的 \(L\) 时，\(G\) 是 \(N\) 与 \(L\) 的内半直积；若选择了同态截面 \(s\)，则
\[
G\cong N\rtimes_\alpha H,\qquad
\alpha_h(n)=s(h)ns(h)^{-1},
\]
同构由 \((n,h)\mapsto ns(h)\) 给出。
\end{theorem}
\begin{proof}
若有同态截面 \(s\)，则由 \(\pi s=\operatorname{id}\) 可知 \(s\) 单射，像 \(L=s(H)\) 是子群，且 \(\pi|_L\) 与 \(s\) 互逆。对任意 \(g\in G\)，令 \(h=\pi(g)\)，则 \(gs(h)^{-1}\in\ker\pi=N\)，所以 \(G=NL\)。若 \(s(h)\in N\)，应用 \(\pi\) 得 \(h=1\)，故 \(N\cap L=\{1\}\)。内半直积判别给出所需同构。

反过来，若 \(\pi|_L\) 是同构，取它的逆再接子群嵌入，得到同态截面。特别地，外半直积投影 \(N\rtimes H\rightarrow H\) 总有同态截面 \(h\mapsto(1,h)\)。
\end{proof}

\subsection{集合截面与同态截面}
\label{b1:subsec:0603:02}

满射 \(\pi:G\rightarrow H\) 的\term[jiemian]{截面}[section]，作为集合映射，仅要求为每个 \(h\) 选择一个原像 \(s(h)\)。有限群时可逐个选择，一般情形可借助选择公理。本节总把单位的代表选为单位，即 \(s(1)=1\)。集合截面存在不等于群扩张分裂，后者要求这个选择还能保持乘法。

\ProseParagraphBreak
考虑加法群的满射
\[
\mathbb Z/4\mathbb Z\longrightarrow\mathbb Z/2\mathbb Z,
\qquad [a]_4\longmapsto[a]_2.
\]
选择 \(s([0]_2)=[0]_4\)、\(s([1]_2)=[1]_4\) 得到集合截面，但
\[
s([1]_2+[1]_2)=[0]_4,\qquad
s([1]_2)+s([1]_2)=[2]_4,
\]
两者不同。另一种可能的非零陪集代表为 \([3]_4\)，其加倍仍为 \([2]_4\)，所以也不能给出同态截面。这里核与商群都同构于 \(C_2\)，中间群却是 \(C_4\)，并非 \(C_2\times C_2\)。

\ProseParagraphBreak
对任意归一化集合截面 \(s\)，每个 \(g\in G\) 仍有唯一表达 \(g=ns(h)\)：先令 \(h=\pi(g)\)，再令 \(n=gs(h)^{-1}\in N\)。因此 \(G\) 的底集合总可用 \(N\times H\) 的坐标描述；困难在于乘法，而非集合的分解。与半直积相比，新出现的项正好记录截面未能保持乘法的差额。

\subsection{扩张数据与半直积}
\label{b1:subsec:0603:03}

给定归一化集合截面，定义
\[
\alpha_h(n)=s(h)ns(h)^{-1},\qquad
c(h,k)=s(h)s(k)s(hk)^{-1}\in N.
\]
每个 \(\alpha_h\) 都是 \(N\) 的自同构，但映射 \(h\mapsto\alpha_h\) 未必是同态。元素 \(c(h,k)\) 称为这个截面的\term[yinzi-ji]{因子集}[factor set]；它等于单位当且仅当截面对相应的一对元素保持乘法。

\begin{proposition}\label{b1:prop:extension-factor-data}
在坐标 \(g=ns(h)\) 下，扩张中的乘法为
\[
(n,h)(m,k)=(n\alpha_h(m)c(h,k),hk).
\]
这些数据满足 \(\alpha_1=\operatorname{id}\)、\(c(1,h)=c(h,1)=1\)，并满足
\begin{align}
\alpha_h\alpha_k
&=\operatorname{Inn}_{c(h,k)}\alpha_{hk},\label{b1:eq:extension-action-defect}\\
c(h,k)c(hk,\ell)
&=\alpha_h(c(k,\ell))c(h,k\ell),\label{b1:eq:extension-factor-identity}
\end{align}
其中 \(\operatorname{Inn}_u(n)=unu^{-1}\) 表示由 \(u\) 给出的\term[nei-zitonggou]{内自同构}[inner automorphism]。
\end{proposition}
\begin{proof}
先把 \(m\) 移过 \(s(h)\)，再用 \(s(h)s(k)=c(h,k)s(hk)\)，得到
\[
ns(h)ms(k)=n\alpha_h(m)c(h,k)s(hk),
\]
这给出乘法公式。由 \(s(h)s(k)=c(h,k)s(hk)\)，两边对 \(N\) 作共轭便得到\eqref{b1:eq:extension-action-defect}。另一方面，将 \(s(h)s(k)s(\ell)\) 分别按左右结合展开，左结合为
\[
c(h,k)c(hk,\ell)s(hk\ell),
\]
右结合为
\[
\alpha_h(c(k,\ell))c(h,k\ell)s(hk\ell).
\]
消去右端共同因子，就得到\eqref{b1:eq:extension-factor-identity}。归一化条件直接来自 \(s(1)=1\)。
\end{proof}
若把截面改为 \(s'(h)=b(h)s(h)\)，其中 \(b(h)\in N\)、\(b(1)=1\)，则直接代入定义给出
\[
\alpha'_h=\operatorname{Inn}_{b(h)}\alpha_h,\qquad
c'(h,k)=b(h)\alpha_h(b(k))c(h,k)b(hk)^{-1}.
\]
计算第二式时，唯一需要移过截面的因子是 \(b(k)\)，由
\(s(h)b(k)=\alpha_h(b(k))s(h)\) 处理。可见因子集依赖截面的选择，同一个扩张可以有不同的 \(c\)。扩张分裂当且仅当能选到使全部 \(c(h,k)=1\) 的截面；这时\eqref{b1:eq:extension-action-defect}说明 \(\alpha\) 成为真正的同态，乘法恢复为半直积。对于非交换的 \(N\)，不能删掉内自同构项并把任意 \(\alpha\) 都当作群作用。
